\subsubsection{The height bound for the terminal evaluation}
\label{subsec:selected-terminal-height}

At the compact end of the chain the height also includes the branching evaluation. Its two interlacing weights determine which coefficients have to be bounded.

Use the reference sequence and the terminal point and row spaces at the refined levels of Section~\ref{subsec:terminal-realization}. At each
real label \(z\), the terminal large and small weights
\(\lambda,\sigma\) have ranks \(N=N_m\) and \(N-1\).
They interlace and have the same last entry. Put
\(A_i=\sigma_i-\lambda_{i+1}\) and \(B_i=\lambda_i-\sigma_i\)
for \(1\le i<N\), and let
\[
 D_{0,z}=\sum_i A_i,\qquad
 n_z=\sum_i\lambda_i-\sum_i\sigma_i=\lambda_1-D_{0,z}.
\]
The harmonic weight on the last diagonal coordinate is
\(n_z=\lambda_N+\sum_i B_i\). In this equality the last
entry of the full weight accounts for the determinant
character in the polynomial model.

\begin{proposition}[Height of the terminal evaluation]\label{prop:terminal-evaluation-height}
For the reference sequence, let \(T\) inside the sum below
denote the selected top labels and \(T_u\) the assigned prime.
The logarithmic squared norm of terminal evaluation in the
indicated point metrics is bounded above by
\[
 -\sum_{z\in T}a_z\log T_u
 +O\!\left(L_a\bigl[|T|b'N+
     f(S_c+m\eta)N+fN^2\bigr]\right).
\]
The implicit constant is uniform for the selected ranges of the parameters.
\end{proposition}

\begin{proof}
\proofstep{Separate the compact branching norm from its Capelli scalar.}
Recall the branching embedding \(J=c_zJ_0\), where
\(J_0h_\sigma=\xi_z\) and \(c_z\) is the nonzero integer
obtained from the capped interlacing intervals. With the
invariant metrics in which both highest vectors have norm
one, \(J_0^*J_0=\|\xi_z\|^2\,1\) by Schur's lemma.
Since \(J=c_zJ_0\), Schur's scalar for the normalized map gives the squared norm of the normalized transpose:
\[
 |c_z|^2\|\xi_z\|^2 .
\] The branching polynomial \(\xi_z\)
is a product of the initial-column flag minors in the
smaller-rank model inside the large one. Its normalized
squared norm is at most one. Indeed realize these minors
as coefficients in tensor powers of exterior powers, with
the wedge basis orthonormal. The tensor of the initial
wedges spans the highest weight of the Cartan summand.

Let \(u\) be the tensor of row wedges, of norm one, and
\(P_{\rm Car}\) the orthogonal projection to the Cartan summand.
The coefficient vector representing \(\xi_z\) is \(P_{\rm Car}u\), so
\[
 \|\xi_z\|^2=\|P_{\rm Car}u\|^2
             =1-\|(I-P_{\rm Car})u\|^2\le1 .
\]
The unit of comparison is the initial wedge tensor
itself, which gives \(h_\lambda\). In particular the
weaker bound \(\log\|\xi_z\|^2\le N(N-1)\log(2a)/2\)
follows since \(a\ge1\) and \(N\ge1\).

For the Capelli scalar, let
\(D_{1,z}=\sum_i\min(\lambda_i-\lambda_{i+1},D_{0,z})\).
This is exactly the number of trial integers in the capped
intervals after omitting the selected \(\sigma_i\).
Each trial contributes \(N-1\) differences. Since all shifted
weights and trial integers lie in the common spread interval,
each difference has absolute value at most \(2a\), by the
selected spread bound. Consequently
\[
             \log |c_z|\le (N-1)D_{1,z}\log(2a).
\]
\label{cor:selected-terminal-height-1}

\proofstep{Use the signed degree in the common class metric.}
At each class the last line is scaled by the fixed word
element \(\pi_u\), with \(\pi_u\overline{\pi_u}=T_u\).
The action in the transported evaluation includes
\(\prod_z z(\pi_u)^{-n_z}\) with the full signed \(n_z\).
Its squared modulus is \(T_u^{-\sum_z n_z}\).
Both metrics in the terminal pairing use the same row
class index and the same reciprocal stabilizer weights.
In particular the large fibre at a row class is the
terminal point fibre prescribed over that class. Applying
Cauchy--Schwarz with these weights gives the squared
pairing bound by
\[
 T_u^{-\sum_z n_z}\prod_z |c_z|^2\|\xi_z\|^2
\]
times the squared norms in the two terminal point metrics.
This estimate starts after the coupled map of
Section~\ref{subsec:selected-norms}; its two-kernel
class sum is already in that map. The evaluation and
quotient-metric comparisons in
Section~\ref{subsec:selected-lift-scalar} give its
bound on the framed spectral row.
\label{cor:selected-terminal-height-2}

\proofstep{Bound the head and the tail by interlacing.}

Write the head Harish--Chandra coordinates at a label as
\(u_0,\ldots,u_{k-1}\), with the interlacing coordinates with smaller weights
\(v_0,\ldots,v_{k-2}\). Write \(b_\alpha,e_\alpha\),
\(0\le\alpha<m\), for the two ends of the successive
tail blocks of size \(t+1\). The coordinates with smaller weights lie
in \([0,S_c]\), \(u_1<v_0\), and
\(|b_\alpha-e_\gamma|\le\eta\). With zero-based indices the head entries in the large and
small weights are, respectively,
\[
 u_i-\frac{k-1}{2}+i-\frac{m(t+1)}2,\qquad
 v_i-\frac{k-2}{2}+i-\frac{m(t+1)}2 .
\]
Their difference on the shared indices is \(u_i-v_i-1/2\).
The subsequent block \(\alpha\) in the large weight is constant at
\(-e_\alpha+k/2+(2\alpha+1-m)(t+1)/2\).

In each tail block the \(t\) later entries agree with
the corresponding entries of the large weight in the
shifted index for \(A_i\). The first entry in this tail block contributes the difference
\[
 e_\alpha-b_\alpha-\tfrac{t+1}{2}
\]
to \(D_{0,z}\).
This follows by substituting the first coordinate
\(-b_\alpha\) and the \(t\) successive coordinates
from \(-e_\alpha\) in the shifted Harish--Chandra row
of the small representation.
\label{cor:selected-terminal-height-3}

The head sum of \(A_i\) is at most
\(v_0-u_{k-1}-k+3/2\) by telescoping with interlacing,
and the tail contributes at most \(m\eta\).
Off the top labels \(u_{k-1}=v_{k-2}-1/2\).
At a top label \(z\), use
\(b_z=-u_{k-1}-(k-1)/2\le2b'\).
It follows that
\[
        D_{0,z}\le S_c+m\eta+2b'\,{\bf1}_T(z),
\]
where \(T\) is the set of top labels.

To bound \(D_{1,z}\),
cap the first gap and the gap between the head and the
first tail block by \(D_{0,z}\). The head gaps left
over telescope from the second to the \(k\)-th entry.
Their sum is bounded by \(D_{0,z}+S_c\): subtract the
single term comparing \(v_{k-2}\) with the last head
entry and use \(u_1<v_0\), \(v_{k-2}\ge0\).
The remaining tail gaps telescope to
\(e_{m-1}-e_0-(m-1)(t+1)\le\eta\), since
\(b_0\le e_0\) and \(|b_0-e_{m-1}|\le\eta\).
Combining the two capped gaps, the remaining head gaps and
the telescoping tail gives, respectively,
\[
 D_{1,z}\le2D_{0,z}+(D_{0,z}+S_c)+\eta
                  =3D_{0,z}+S_c+\eta .
\]

At a top label \(\lambda_1=a_z-(N-k)/2\).
Otherwise \(u_0-v_0\ge S_c\) and \(v_0\ge0\)
give \(\lambda_1\ge S_c-(N-1)/2\), which is sufficient
even when \(n_z\) is negative. Therefore
\[
 \sum_z n_z\ge \sum_{z\in T}a_z-2|T|b'-f(S_c+m\eta+N).
\]
Sum the logarithmic bounds for the branching scalars.
Using \(\log(2a)=O(L_a)\) and
\(\log T_u\le25mL\le25N L_a\) gives the logarithmic
squared evaluation price
\[
 -\sum_{z\in T}a_z\log T_u
 +O\!\left(L_a\bigl[|T|b'N
        + f(S_c+m\eta)N+fN^2\bigr]\right).
\]
This is \(h_A\) with an absolute uniform error constant,
on the same reference sequence as the \(h_B(j)\) bounds.
\label{cor:selected-terminal-height-4}
\end{proof}

